\chapter{Sintonización PID mediante Particle Swarm Optimization (PSO)}
\label{ch:pso}

\section{Introducción del capítulo}
El algoritmo Particle Swarm Optimization es el eje metodológico central del borrador base en
el que se apoya este libro. Su interés radica en que ofrece una búsqueda cooperativa simple,
de bajo costo computacional relativo y con muy buenos resultados en problemas de
sintonización PID donde la evaluación de cada solución depende de una simulación dinámica.
En el caso de la microred de Limbani, el PSO destaca por su capacidad para adaptarse al
comportamiento no lineal del conjunto gobernador--turbina.

\section{Objetivo del capítulo}
Este capítulo explica la lógica del PSO, resume su formulación matemática y muestra cómo se
adapta al ajuste de las ganancias $K_p$, $K_i$ y $K_d$. El objetivo es dejar claro por qué este
método constituye una alternativa sólida frente a la sintonización clásica cuando se trabaja con
criterios como ITAE y escenarios de perturbación severa.

\section{Inspiración del algoritmo}
PSO se inspira en el comportamiento colectivo de bandadas y enjambres, donde cada
individuo modifica su movimiento usando tanto su experiencia propia como la información del
grupo. Trasladado al lenguaje de optimización, cada partícula representa una solución
candidata y navega por el espacio de búsqueda recordando su mejor posición histórica y la
mejor posición encontrada por el enjambre. Esa memoria dual permite combinar exploración
global con convergencia rápida hacia regiones prometedoras.

\section{Modelo matemático del PSO}
Cada partícula $i$ posee una posición $x_i^k$ y una velocidad $v_i^k$ en la iteración $k$. La
actualización típica se expresa como
\[
v_i^{k+1}=w v_i^k + c_1 r_1\left(pbest_i-x_i^k\right)+c_2 r_2\left(gbest-x_i^k\right),
\]
\[
x_i^{k+1}=x_i^k+v_i^{k+1}.
\]
Aquí, $w$ es el factor de inercia, $c_1$ y $c_2$ son coeficientes de aceleración, y $r_1$, $r_2$
son valores aleatorios uniformes. En la aplicación de este libro, la posición de cada partícula
corresponde al vector de ganancias $\left[K_p,K_i,K_d\right]$ y su calidad se mide mediante el
valor del ITAE obtenido en simulación.

\section{Pseudocódigo}
\begin{algorithm}[H]
\caption{Esquema general del PSO para sintonización PID}
\begin{algorithmic}[1]
\State Inicializar una población de partículas en el espacio $\left[K_p,K_i,K_d\right]$
\State Evaluar cada partícula mediante la simulación de la planta y el cálculo de ITAE
\State Guardar $pbest_i$ y $gbest$
\While{no se cumpla el criterio de parada}
    \For{cada partícula}
        \State Actualizar velocidad y posición
        \State Limitar las ganancias al rango admisible
        \State Recalcular el costo de la partícula
        \State Actualizar $pbest_i$ y $gbest$ si mejora el desempeño
    \EndFor
\EndWhile
\State Devolver la mejor tripleta $\left[K_p,K_i,K_d\right]$
\end{algorithmic}
\end{algorithm}

\section{Parámetros de ajuste}
Los parámetros del PSO condicionan fuertemente su desempeño. Un valor intermedio de inercia
favorece la búsqueda global sin impedir la convergencia; los coeficientes cognitivo y social
equilibran la atracción hacia la experiencia individual y la colectiva; y el tamaño del enjambre
determina el nivel de diversidad exploratoria. En el borrador de referencia se consideran
configuraciones del orden de 15 a 20 partículas y un número moderado de iteraciones, lo que
resulta suficiente para obtener mejoras claras sin disparar el costo computacional.

\section{Aplicación a la sintonización PID}
La aplicación del PSO al problema de sintonización es directa: cada partícula codifica una
tripleta de ganancias del PID y se evalúa haciendo correr el modelo de la microred ante una o
varias perturbaciones. La función objetivo prioriza respuestas con bajo error acumulado,
mínimas oscilaciones y recuperación rápida de frecuencia. Esta formulación permite que la
búsqueda se guíe por el comportamiento real del sistema, en lugar de depender solo de reglas
empíricas derivadas de una condición crítica idealizada.

\section{Cierre del capítulo}
PSO representa un punto de equilibrio muy atractivo entre simplicidad algorítmica y calidad de
la sintonización obtenida. Su desempeño justifica que sea tomado como referencia principal en
la comparación. Aun así, no es la única alternativa disponible; el capítulo siguiente introduce
el Grey Wolf Optimizer como otra estrategia de búsqueda basada en inteligencia colectiva.
